\section{Related Work}

% \subsection{Visual Reinforcement Learning}
\paragraph{Visual Reinforcement Learning} Successes of visual representation learning in computer vision~\citep{vincent2008extracting, doersch2015unsupervised,Wang_UnsupICCV2015, noroozi2016unsupervised,zhang2017split,gidaris2018unsupervised} has inspired successes in visual RL, where coherent representations are learned alongside RL. Works such as SAC-AE~\citep{yarats2019improving}, PlaNet~\citep{hafner2018planet}, and SLAC~\citep{lee2019stochastic}, demonstrated how auto-encoders~\citep{finn2015deepspatialae} could improve visual RL. Following this, other self-supervised objectives such as contrastive learning in CURL~\citep{srinivas2020curl} and ATC~\citep{stooke2020decoupling}, self-prediction in SPR~\citep{schwarzer2020data}, contrastive cluster assignment in Proto-RL~\citep{yarats2021proto}, and augmented data in DrQ~\citep{yarats2021image} and RAD~\citep{laskin2020reinforcement}, have significantly bridged the gap between state-based and image-based RL. Future prediction objectives~\citep{hafner2018planet,hafner2019dream,yan2020learning,finn2015deepspatialae,pinto2016curious,agrawal2016learning} and other auxiliary objectives~\citep{jaderberg2016reinforcement,zhan2020framework,young2020visual,chen2020robust} have shown improvements on a variety of problems ranging from gameplay, continuous control, and robotics. In the context of visual control settings, clever use of augmented data~\citep{yarats2021image,laskin2020reinforcement} currently produces state-of-the-art results on visual tasks from DMC~\citep{tassa2018dmcontrol}. %This paper improves on DrQ~\citep{yarats2021image} by introducing several key modifications, and yields an algorithm 

\paragraph{Humanoid Control}  The humanoid control problem first presented in~\cite{tassa2012trajopt}, has been studied as one of the hardest control problems due to its large state and action spaces. The earliest solutions to this problem use ideas in model-based optimal control to generate policies given an accurate model of the humanoid~\cite{}. Subsequent works in RL have shown that model-free policies can solve the humanoid control problem given access to proprioceptive state observations. However, solving such a problem from visual observations has been a challenging problem, with leading RL algorithms making little progress to solve the task~\citep{tassa2018dmcontrol}. Recently, \cite{hafner2020dreamerv2} was able to solve this problem through a model-based technique in around 30M environment steps and $340$ hours of training on a single GPU machine. \drqs{}, presented in this paper, marks the first model-free RL method that can solve humanoid control from visual observations, taking also around 30M steps and $86$ hours of training on the same hardware.

%\paragraph{Democratizing Large-scale RL}\todo{this seems out of place under related works. Maybe it can be moved to a different section?}
%Significant efforts have been devoted to large-scale parallel RL training infrastructure \cite{}. These have been used to demonstrate impressive applications of RL to Go, StarCraft, etc. \cite{}. However, their computational requirements are beyond all but a few selected research labs. \drqs{} has the attractive property that it is highly efficient, being practical to train on a single GPU, within hours rather than days for most tasks. 
% by What is this environment? MPC/optimal control based papers. State-based papers on this environment with rough sample complexity estimate for SOTA state-based. Dreamer as the first vision RL method. But it is model-based, which means complex + takes flops to run. This paper we show that humanoid control from pixels can be solved through model-free learning.

% In contrast to computer vision, data augmentation is rarely used in RL. Certain approaches implicitly adopt it, for example~\citet{levine2018learning,kalashnikov2018qt} use image augmentation as part of the AlexNet training pipeline without analysing the benefits occurring from it, thus being overlooked in subsequent work. HER~\citep{andrychowicz2017hindsight} exploits information about the observation space by goal and reward relabeling, which can be viewed as a way to perform data augmentation. Other work uses data augmentation to improve generalization in domain transfer ~\citep{cobbe2018quantifying}.  However, the classical image transformations used in vision have not previously been shown to definitively help on standard RL benchmarks. Concurrent with our work, RAD~\citep{laskin2020reinforcement} performs an exploration of different data augmentation approaches, but is limited to transformations of the image alone, without the additional augmentation of the Q-function used in our approach. Moreover, RAD can be regarded as a special case of our algorithm. Multiple follow ups to our initial preprint appeared on ArXiv~\citep{raileanu2020automatic,okada2020dreaming}, using similar techniques on other tasks, thus supporting the effectiveness and generality of data augmentation in RL. 



% \textbf{Continuous Control from Pixels} There are a variety of methods addressing the sample-efficiency of RL algorithms that directly learn from pixels. The most prominent approaches for this can be classified into two groups, model-based and model-free methods. The model-based methods attempt to learn the system dynamics in order to acquire a compact latent representation of high-dimensional observations to later perform policy search~\citep{hafner2018planet,lee2019slac,hafner2019dream}. In contrast, the model-free methods either learn the latent representation indirectly by optimizing the RL objective~\citep{barth-maron2018d4pg,abdolmaleki2018maximum} or by employing auxiliary losses that provide additional supervision~\citep{yarats2019improving,srinivas2020curl,sermanet2018time,dwibedi2018visrepr}. Our approach is complementary to these methods and can be combined with them to improve performance.


% \subsection{Humanoid}

